\begin{figure*}
\begin{minipage}[b]{0.25\textwidth}
{\small\begin{verbatim}
unsigned
g(unsigned a) {
  switch (a % 4) {
  case 0:
    a += 3;
    break;
  case 1:
    a += 2;
    break;
  case 2:
    a += 1;
    break;
  }
  return a & 3;
}
\end{verbatim}}
\centering (a)
\end{minipage}
%
\begin{minipage}[b]{0.45\textwidth}
{\small\begin{verbatim}
define i32 @g(i32 %0) {
label %1:
  %2 = urem i32 %0, 4
  switch i32 %2, label %9 [
    i32 0, label %3
    i32 1, label %5
    i32 2, label %7
  ]

label %3:
  %4 = add i32 %0, 3
  br label %9

label %5:
  %6 = add i32 %0, 2
  br label %9

label %7:
  %8 = add i32 %0, 1
  br label %9

label %9:
  %.0 = phi i32 [ %0, %1 ], [ %8, %7 ],
                [ %6, %5 ], [ %4, %3 ]
  %10 = and i32 %.0, 3
  ret i32 %10
}
\end{verbatim}}
\centering (b)
\end{minipage}
%
\begin{minipage}[b]{0.33\textwidth}
{\small\begin{verbatim}
%0 = block 4
%1:i32 = var
%2:i32 = urem %1, 4:i32
%3:i1 = ne 0:i32, %2
%4:i1 = ne 1:i32, %2
%5:i1 = ne 2:i32, %2
blockpc %0 0 %3 1:i1
blockpc %0 0 %4 1:i1
blockpc %0 0 %5 1:i1
blockpc %0 1 %2 2:i32
blockpc %0 2 %2 1:i32
blockpc %0 3 %2 0:i32
%6:i32 = add 1:i32, %1
%7:i32 = add 2:i32, %1
%8:i32 = add 3:i32, %1
%9:i32 = phi %0, %1, %6, %7, %8
%10:i32 = and 3:i32, %9
infer %10
\end{verbatim}}
$\Rightarrow$
{\small\begin{verbatim}
result 3:i32
\end{verbatim}}
\centering (c)
\end{minipage}
\caption{A C or C++ function that can be optimized to ``\texttt{return 3}''
  (a), its representation in LLVM~IR (b), and the corresponding \tool{}
  LHS and RHS
  (c). The \tool{} IR has no control
  flow, but rather uses \texttt{blockpc} instructions to represent
  dataflow facts derived from converging control flow. These are
  crucial in giving \tool{} the information that it needs to
  synthesize the optimization.}
\label{fig:blockpc}
\end{figure*}
